\documentclass[10pt,a4paper]{article}
\usepackage{iftex}
\ifPDFTeX
  \usepackage[T1]{fontenc}
  \usepackage{lmodern}
  \usepackage[utf8]{inputenc}
\else
  \usepackage{fontspec}        % XeTeX/LuaTeX (tectonic) : UTF-8 natif
\fi
\usepackage[french]{babel}
\usepackage[a4paper,margin=1.9cm]{geometry}
\usepackage{amsmath,amssymb}
\usepackage{booktabs}
\usepackage{array}
\usepackage{enumitem}
\usepackage{xcolor}
\usepackage{hyperref}
\hypersetup{colorlinks=true,linkcolor=blue!50!black,urlcolor=blue!50!black}
\usepackage{titlesec}
\titlespacing{\section}{0pt}{0.9em}{0.3em}
\titleformat{\section}{\large\bfseries}{\thesection.}{0.5em}{}
\setlength{\parindent}{0pt}
\setlength{\parskip}{0.35em}
\setlist[itemize]{leftmargin=1.3em,itemsep=1pt,topsep=2pt}
\newcommand{\code}[1]{\texttt{\small #1}}

\title{\vspace{-1.2em}\textbf{Du corpus brut à la table de backtest}\\
\large Nettoyage, validations et garanties --- fiche de synthèse (2014--2026)}
\author{}
\date{}

\begin{document}
\maketitle
\vspace{-2.2em}

% ══════════════════════════════════════════════════════════════════
\section{D'où on part : un corpus complet, identifié, vérifié}

\textbf{La collecte : à la source officielle, pas chez un agrégateur.}
\begin{itemize}
\item STOCK Act (2012) : chaque élu déclare ses transactions sous 45~jours dans un \emph{PTR}
      (\emph{Periodic Transaction Report}).
\item Deux dépôts officiels : l'\textbf{index annuel du Clerk} (House --- liste chaque PTR déposé =
      univers de référence re-vérifiable) et l'\textbf{eFD du Sénat}.
\item Chaque document est téléchargé puis lu : extraction directe (dépôts électroniques) ou OCR Vision
      (PDF scannés).
\item Toute la chaîne est contrôlée : chaque ligne garde l'identifiant de son document source.
\end{itemize}

\begin{center}\small
\begin{tabular}{p{7.6cm}p{8.6cm}}
\toprule
Transactions uniques (2 chambres, 2014--2026) & \textbf{169\,000} \quad (House 151\,989 · Sénat 17\,011) \\
\quad\emph{dont} : brutes 170\,920, re-divulgations dédupliquées & 1\,920 \\
Index officiel House Clerk : PTR traités & \textbf{100\,\% des 8\,252} (7\,568 avec transactions · 582 manuscrits gated · 102 vides documentés) \\
Identité rattachée (bioguide) & \textbf{100\,\%} (367 House + 78 Sénat) \\
Dates cohérentes / délai médian de publication & 99{,}8\,\% / 27 jours \\
Montants renseignés & 99{,}4\,\% \\
\bottomrule
\end{tabular}
\end{center}

\emph{Lecture du tableau} : « re-divulgations » = trade republié dans un amendement (dédoublonné par clé
naturelle) · « manuscrits gated » = PTR remplis à la main, écartés (motif au \S4) · « vides documentés »
= dépôts sans transaction, justifiés un par un.

\textbf{La validation : trois sources tierces} (\emph{Quiver Quant} + deux projets open-source), en
validation seulement --- rien n'y est copié. Trois questions :

\textbf{1) A-t-on tout ce qu'eux ont ?}
\begin{itemize}
\item Quiver, \emph{sans} la date : \textbf{93{,}5\,\%} (House) / \textbf{92{,}1\,\%} (Sénat) de ses
      combinaisons retrouvées chez nous. \emph{(Combinaison = ce déposant a déclaré au moins un trade sur
      ce titre dans ce sens ; on compare d'abord sans la date, pour tolérer les conventions de date.)}
\item Quiver, \emph{avec} la date : \textbf{70\,265} (House) + \textbf{8\,766} (Sénat) trades appariés
      \textbf{au jour exact}, +\,845 à $\leq$\,10 jours.
\item Les seuls candidats sérieux d'« erreur de date » : 545 paires issues d'une même déclaration
      (écart typique 11~j) --- vérification faite, c'est une convention de date de Quiver ; le vrai
      contrôle des dates est l'audit des PDF officiels.
\item Les combinaisons de Quiver qu'on ne retrouve \emph{pas} (1\,704 côté House, 360 côté Sénat) ne sont
      presque jamais de vrais manques. En regardant de près :
      \begin{itemize}
      \item la plupart (1\,578 House, 7 Sénat) sont bien chez nous, mais dans des PTR manuscrits qu'on a
            écartés \emph{exprès} (voir \S4) ;
      \item d'autres (95 House, 353 Sénat) ne sont pas backtestables --- actif non coté ou ticker
            fragmenté même chez Quiver ;
      \item il ne reste, au pire, que \textbf{31 combinaisons cotées côté House (0 côté Sénat)} qui
            pourraient être de vrais manques.
      \end{itemize}
\item \textbf{senate-stock-watcher} (base publique et indépendante du Sénat, 2014--2020, 6\,231 trades
      avec ticker) : on cherche chacune de \emph{ses} lignes chez nous → \textbf{100\,\%} retrouvées. Ses
      36 exceptions ne sont pas des manques : sur un échange de titres, il écrit 2 lignes là où nous en
      écrivons 1.
\item \textbf{Mirror house-stock-watcher} (une seconde lecture indépendante des \emph{mêmes} PDF de la
      Chambre, 23\,568 trades) : \textbf{99{,}5\,\%} de ses lignes 2018--2026 sont chez nous. Avant 2018,
      ses rares écarts nous ont \emph{servi} : chacun pointait un trade qu'on avait manqué --- tous
      récupérés et corrigés depuis.
\end{itemize}

\textbf{2) A-t-on \emph{plus} qu'eux ?}
\begin{itemize}
\item \textbf{+23\,535} trades cotés nets côté House, \textbf{+2\,951} côté Sénat ; au trade daté :
      39\,207 + 3\,140 trades chez nous absents de Quiver.
\item Le surplus est réel, pas un artefact : \textbf{Quiver est mince avant 2017} (il ne corrobore que
      16{,}6--32{,}6\,\% de nos actions 2014--2016).
\item Nos lignes « nous-seul » portent des tickers vérifiés sur les PDF officiels.
\end{itemize}

\textbf{3) Sur les trades communs, dit-on la même chose ?}
\begin{itemize}
\item $\sim$96\,000 paires appariées 1-à-1 (membre, ticker, date).
\item Accord sur le \textbf{sens} : 96{,}8\,\% (House) / 99{,}9\,\% (Sénat) ; sur le \textbf{montant} :
      94{,}2\,\% / 99{,}8\,\%.
\item Désaccords échantillonnés vérifiés sur les PDF officiels : notre transcription est fidèle au document.
\end{itemize}

% ══════════════════════════════════════════════════════════════════
\section{L'entonnoir de nettoyage : 4 retraits, tous justifiés}

\emph{Philosophie (écrite en tête du notebook)} : \textbf{on ne retire que l'avéré inutilisable ; tout le
doute est gardé et flagué} --- aucune ligne écartée en silence, la décision revient à la recherche, pas au
nettoyage.

\begin{center}\small
\begin{tabular}{clrrr}
\toprule
Étape & Règle (pourquoi) & retirées & restantes & \% \\
\midrule
--- & Corpus unique (\code{load\_final}, réparations incluses) & --- & 169\,000 & 100{,}0 \\
A & Dates présentes \& cohérentes (impossible de dater l'entrée sinon) & 3\,252 & 165\,748 & 98{,}1 \\
B & Actions + ETF \textbf{cotés} (un backtest exige un prix de marché) & 29\,771 & 135\,977 & 80{,}5 \\
C & Direction claire achat/vente (les « échanges » n'ont pas de sens directionnel) & 826 & 135\,151 & 80{,}0 \\
D & Montant présent (dimensionner une position exige un notionnel) & 687 & \textbf{134\,464} & \textbf{79{,}6} \\
\bottomrule
\end{tabular}
\end{center}

\textbf{Étape A --- les trois critères de retrait :}
\begin{itemize}
\item dates illisibles ou absentes (impossible de dater l'entrée) ;
\item publication \emph{avant} la transaction --- impossible, coquille du déposant ;
\item année implausible : transaction avant 2012 (avant le STOCK Act) ou après l'année du dépôt ;
\item en revanche, les \emph{dépôts tardifs} (12{,}2\,\% $>45$~j) sont GARDÉS et flagués : information
      réelle, jouable à sa date de publication.
\end{itemize}

\textbf{Étape B --- ce qu'on \emph{retire} (29\,771 lignes) :}
\begin{itemize}
\item 27\,006 lignes \emph{sans aucun symbole} (munis, obligations, sociétés privées --- pas de prix possible) ;
\item 164 tickers malformés (CUSIP, fragments OCR) ;
\item 2\,601 familles non cotées (options, bons du Trésor).
\end{itemize}

\textbf{Étape B --- ce qu'on \emph{garde} grâce au ticker-first :}
\begin{itemize}
\item \textbf{21\,878 lignes GARDÉES} : leur champ « type d'actif » est vide, mais leur ticker est une
      action cotée valide → on les garde (on décide sur le ticker, pas sur le champ déclaré).
\item pourquoi ça compte : ce champ « type d'actif » est vide dans 16 à 31\,\% des cas ; le filtrer
      naïvement (comme dans une première version) jetait à tort ces 21\,878 vraies actions --- bug
      trouvé, corrigé, nettoyage centralisé ici.
\end{itemize}

\textbf{Ce que la table contient :}
\begin{itemize}
\item chambres : House 91{,}3\,\% / Sénat 8{,}7\,\% ; sens : achats 68\,250 / ventes 66\,214 ;
\item classes : actions 128\,974 (secteur GICS 100\,\%) · ETF diversifiés 2\,022 · ETF sectoriels 168 ·
      3\,300 cotés à classe non résolue (flagués) ;
\item les deux ères couvertes : House 56\,706 en 2014--2019, 66\,108 en 2020--2026.
\end{itemize}

% ══════════════════════════════════════════════════════════════════
\section{Les garanties : rien ne bouge sans qu'on le voie}

\textbf{Le doute est gardé et flagué} (jamais retiré en silence) --- la recherche a toute l'info pour le
traiter :
\begin{itemize}
\item \code{flag\_late\_filing} --- publication au-delà des 45~jours légaux (16\,349 lignes, 12{,}2\,\%) :
      c'est de l'information réelle, jouable à sa date de publication ; on la garde.
\item \code{flag\_very\_late\_filing} --- publication au-delà d'un an (3\,075, 2{,}3\,\%) : amendements ou
      régularisations très tardives ; la recherche peut les filtrer si elle veut.
\item \code{is\_delisted} --- le titre a quitté la Bourse pendant la période (3\,542, 2{,}6\,\%), typé
      \emph{faillite} vs \emph{rachat} : le prix de sortie n'est pas le même (\textasciitilde 0 en
      faillite, prix du deal en rachat) --- le type indique à la recherche quelle valeur finale prendre.
\item \textbf{lots multi-comptes} (11\,873, 8{,}8\,\%) : le même trade déclaré sur plusieurs comptes
      (Self / conjoint / joint) --- des positions \emph{réelles distinctes}, pas des doublons ; on les
      repère par \code{owner}/\code{occurrence\_index}/\code{lot\_size} → \textbf{ne jamais dédoublonner}.
\item \textbf{Assertions finales sur 100\,\% des 134\,464 lignes} : bioguide, ticker, montant, direction
      $\in\{$buy, sell$\}$, chronologie (publication $\geq$ transaction), clé naturelle --- le notebook
      refuse de publier la table sinon.
\end{itemize}

% ══════════════════════════════════════════════════════════════════
\section{Les limites, écrites \emph{avant} la recherche}

\begin{itemize}
\item \textbf{582 PTR manuscrits écartés} par décision de périmètre : la corroboration externe
      s'effondre sur les manuscrits (12{,}9\,\%), donc invérifiables contre une vérité-terrain --- choix
      conservateur, documenté document par document ; 7{,}1\,\% de l'index, concentrés avant 2020.
\item \textbf{Couverture prix 87{,}8\,\%} des trades (74\,\% en 2014 $\to$ 99\,\% en 2026) : Yahoo ne
      sert plus certains délistés sans successeur $\Rightarrow$ léger biais de survie.
\end{itemize}

\end{document}

